% =============================================================
%  Section 4 -- Data
%  Label: sec:data
% =============================================================

\section{Data}\label{sec:data}

\subsection{China Family Panel Studies}\label{sec:cfps}

The China Family Panel Studies (CFPS) is a nationally
representative biennial household panel survey administered
by the Institute of Social Science Survey (ISSS) at Peking
University \citep{CFPSData}.
The baseline wave in 2010 drew a stratified, multi-stage
probability sample of approximately 16,000 households in
25 mainland provinces, covering roughly 95 percent of the
national population.
The survey has been fielded in 2010, 2012, 2014, 2016, 2018,
2020, and 2022.
I use the four waves from 2016 to 2022, which overlap with
the availability of province-level meat CPI data from the NBS.
For auxiliary cross-sectional checks in companion specifications,
I draw CHFS source files from the China Household Finance Survey
and Research Center at Southwestern University of Finance and
Economics \citep{CHFSData}.

Response rates are high by international panel standards.
The baseline household response rate was approximately
81 percent, with re-contact rates between 75 and 85 percent
in subsequent waves.
Attrition is non-random---younger, more educated, and
urban-resident households are more likely to exit---and I
include age, education, and urban indicators as controls
where available.

The inflation expectation variable is a three-point ordinal
question (item qp201) that asks respondents: ``Compared with
the current period, do you think the price level in the coming
year will rise, stay the same, or fall?''
I code the response as $\mu_{it} \in \{-1, 0, 1\}$, where
$1$ indicates expected price increases and $-1$ indicates
expected decreases.
The question wording refers to ``price level'' (物价水平)
rather than the CPI specifically, which is consistent with
the salience channel: if respondents map ``price level'' to
the prices of goods they buy frequently, their responses will
be disproportionately influenced by salient price changes.

The coarseness of the three-point scale attenuates the estimated
coefficients: respondents who revise their beliefs upward by
different amounts all register as $+1$.
This attenuation bias works against finding overreaction, so
any detected overreaction is a lower bound.

The estimation sample contains 90,133 household-wave
observations across 31 provinces and four waves.
The forecast-error sample is smaller (69,529 observations)
because the 2022 wave lacks forward realised headline CPI
at the time of data construction.
Demographic controls---age, a high-education indicator, and
an urban-residence indicator---are available for a subsample
of approximately 22,100 observations.

\subsection{Province-Level CPI Data}\label{sec:nbs_cpi}

The National Bureau of Statistics (NBS) publishes monthly
consumer price indices at the province level for major
sub-categories.
I use the meat and poultry sub-index, the grain sub-index,
and the headline CPI index.
Province-level meat CPI is available from January 2016 onward;
grain and headline CPI are available from January 2012.

Geographic coverage spans all 31 mainland provinces.
Two potential measurement concerns arise.
The outlet sample may not fully represent rural areas,
where wet markets are the dominant channel for fresh protein.
If rural meat prices move differently from the sampled outlets,
the NBS meat sub-index may understate the price variation
experienced by rural CFPS households, attenuating the
estimated coefficients toward zero.
The transition from separate urban and rural CPI
baskets to a unified basket in 2016 may also create a level shift
in some province-level series.
I avoid this problem by beginning the analysis in 2016 and
using cumulative log changes, which difference out level shifts.

The NBS reports these sub-indices in \emph{previous month
$= 100$} form: a value of 105 means the price level rose
5 percent relative to the preceding month.
I convert each monthly observation to a log change by
dividing by 100 and taking the natural logarithm; because the
base is the previous month, this yields the exact
month-on-month log growth rate.
I then sum these monthly log changes over the inter-wave
period to obtain the cumulative treatment.

\subsection{Treatment Construction}\label{sec:shock_construction}

The treatment variable is the cumulative log change in
provincial meat CPI between two consecutive CFPS waves:
\begin{equation}\label{eq:meat_shock_data}
  \mathrm{MeatShock}_{p,\,t-1 \to t}
  = \sum_{m \in \mathcal{M}(t-1,t)}
      \log\!\left(\frac{\mathrm{MeatCPI}_{pm}}{100}\right),
\end{equation}
where $\mathrm{MeatCPI}_{pm}$ is the NBS meat and poultry
sub-index for province $p$ in month $m$, and
$\mathcal{M}(t-1,t)$ is the set of months between the
midpoints of two consecutive CFPS fieldwork periods.
The grain shock is constructed identically from the
grain sub-index.

I define the inter-wave period using the midpoint of each
wave's fieldwork window.
The 2016 wave was fielded primarily between June and December
2016; the 2018 wave between June and November 2018; the 2020
wave between July 2020 and January 2021 (delayed by
COVID-19); and the 2022 wave between June and December 2022.
This midpoint convention ensures that the treatment
captures the price changes that the typical respondent
experienced between consecutive interview dates.

The identifying variation operates at the province$\times$wave
level.
The sample contains 93 province$\times$wave cells across
31 provinces and three to four waves, depending on forward
CPI availability.
All standard errors are clustered at the province level
(31 clusters), and inference uses the wild cluster bootstrap
to account for the small number of clusters.

\subsection{Descriptive Statistics}\label{sec:desc_stats}

\begin{table}[!htbp]
\centering
\caption{Descriptive Statistics}
\label{tab:summary_stats}
\begin{threeparttable}
\begin{tabular}{lrrrrrr}
\toprule
Variable & $N$ & Mean & SD & Min & Median & Max \\
\midrule
Inflation expectation ($\mu_{it}$) & 90,133 & -0.32 & 0.81 & -1.00 & -1.00 & 1.00 \\
Forecast error ($FE_{it}$) & 69,533 & 4.67 & 1.88 & -0.21 & 5.07 & 8.14 \\
Meat shock ($s_{pt}$) & 90,133 & 0.05 & 0.28 & -0.40 & -0.01 & 0.57 \\
Grain shock & 90,133 & 0.04 & 0.02 & -0.07 & 0.03 & 0.12 \\
Realized headline CPI (\%, ann.) & 69,533 & 4.03 & 0.96 & 1.79 & 3.86 & 6.14 \\
Age & 22,320 & 40.19 & 14.27 & 18.00 & 39.00 & 70.00 \\
High education (indicator) & 90,133 & 0.30 & 0.46 & 0.00 & 0.00 & 1.00 \\
Urban (indicator) & 87,414 & 0.67 & 0.47 & 0.00 & 1.00 & 1.00 \\
\bottomrule
\end{tabular}
\begin{tablenotes}[flushleft]
\footnotesize
\item \textit{Note.} Sample restricted to observations with non-missing province code, inflation expectation, and meat shock.
\item Inflation expectation coded as $\{-1, 0, 1\}$ (prices fall, stable, rise).
\item Forecast error = realized headline CPI (\%) $-$ expectation scaled to CPI units ($\{-2, 0, 2\}$).
\item Demographics available for a subsample; $N$ varies across rows.
\end{tablenotes}
\end{threeparttable}
\end{table}


Table~\ref{tab:summary_stats} reports summary statistics for
the estimation sample.

The mean inflation expectation is $-0.32$ on the $[-1, 1]$
scale, indicating that the majority of respondents expect
prices to stay the same or fall.
This finding is consistent with two features of the data.
The sample period includes the post-ASF price reversal of
2021--2022, during which headline CPI inflation fell below
1 percent. The $\{-1, 0, 1\}$ coding also treats ``stay the
same'' as the midpoint.

The mean forecast error is $4.67$ percentage points.
The positive sign reflects the mechanical asymmetry of the
ordinal coding: when respondents report expectations of zero
or below but realised CPI inflation is positive, the forecast
error is positive.
The regression analysis controls for this level effect through
fixed effects; the coefficient of interest captures the
\textit{differential} forecast error in high-shock versus
low-shock provinces, not the mean level.

The mean meat shock is $0.05$ log points, with a standard
deviation of $0.28$.
The near-zero mean reflects offsetting waves: the 2018--2020
pair captures the ASF peak (large positive shocks), while the
2020--2022 pair captures the price reversal (large negative
shocks).
A one-standard-deviation shock of $0.28$ log points
corresponds to a cumulative meat-price increase of roughly
32 percent ($e^{0.28} - 1 \approx 0.32$) over the inter-wave
period.

Forward realised headline CPI averages $4.03$ percent
(annualised), with moderate cross-province variation driven
by the elevated food prices during the 2018--2020 inter-wave
period.
